% --------------------------------------------------
% Language and typography — configuration for pdfLaTeX
% --------------------------------------------------

% Encodage des fontes utilisées pour composer le texte.
%
% T1, aussi appelé encodage Cork, définit 256 positions de caractères
% dans une fonte, notamment des lettres accentuées européennes.
% Ce n'est ni une police particulière, ni l'encodage du fichier source.
%
% Avec l'ancien encodage OT1, certains caractères accentués sont
% construits en superposant un accent et une lettre.
% T1 permet de les utiliser comme des caractères à part entière,
% ce qui facilite notamment la césure des mots accentués.
%
% Exemple : écrire "café" ou le nom d'un auteur avec des accents.
% T1 reste pertinent dans un document en anglais contenant ces mots.
%
% Avec LuaLaTeX ou XeLaTeX, utiliser plutôt fontspec :
% ne pas reprendre cette configuration fontenc/inputenc.
\usepackage[T1]{fontenc}

% Règles de césure et intitulés automatiques en anglais.
% Exemple : \tableofcontents porte le titre "Contents".
% Influence également le format de \today.
\usepackage[french]{babel}

% Encodage du FICHIER SOURCE : indique comment lire les caractères
% que tu saisis, par exemple "é".
%
% À distinguer de fontenc :
% - inputenc : lecture des caractères du fichier .tex ;
% - fontenc  : encodage des fontes pour composer le document.
%
% UTF-8 est déjà le défaut des versions récentes de pdfLaTeX :
% cette ligne est donc facultative.
\usepackage[utf8]{inputenc}

% Utilise les fontes vectorielles Latin Modern, une version enrichie
% des fontes Computer Modern traditionnelles de LaTeX.
% Elles restent nettes à tout niveau de zoom et permettent à microtype
% d'ajuster légèrement la largeur des caractères (expansion).
% À charger après \usepackage[T1]{fontenc}.
\usepackage{lmodern}

% Améliore automatiquement la typographie, notamment grâce à de
% légers ajustements de largeur des caractères et à leur placement
% en bord de ligne, selon le moteur et les fontes.
% Aucune commande particulière n'est nécessaire dans le texte.
\usepackage{microtype}

% Désactive la césure automatique des mots.
% Attention : peut entraîner des espaces trop grands entre les mots
% ou des lignes qui débordent lorsque le texte est justifié.
\usepackage[none]{hyphenat}

% Supprime les alinéas et ajoute un espace entre les paragraphes.
\usepackage{parskip}

% Personnalise la mise en forme des listes.
\usepackage{enumitem}

% Aligne les tirets sur la marge gauche.
\setlist[itemize]{
    label=--,
    labelindent=0pt,
    leftmargin=*,
    align=left
}

% --------------------------------------------------
% Mathematics and figures
% --------------------------------------------------

% Outils de composition mathématique : équations alignées,
% matrices, texte dans les formules et références aux équations.
%
% Exemple :
% \begin{align}
%   u &= x + y, \\
%   v &= x - y.
% \end{align}
%
% Exemple de texte dans une formule :
% \[ u = 0 \quad \text{on the boundary}. \]
%
% Exemple de référence : Equation~\eqref{eq:continuity}.
\usepackage{amsmath}

% Symboles mathématiques supplémentaires.
% Charge aussi amsfonts, qui fournit notamment \mathbb.
%
% Exemple : $\mathbf{u} \in \mathbb{R}^{2}$.
% Exemple : $A \subsetneq B$.
\usepackage{amssymb}

% Dessins vectoriels réalisés directement dans le document :
% grilles, axes, flèches et schémas.
%
% Exemple :
% \begin{tikzpicture}
%   \draw[step=1cm, gray] (0,0) grid (3,2);
%   \draw[->] (0,0) -- (3.5,0) node[right] {$x$};
%   \draw[->] (0,0) -- (0,2.5) node[above] {$y$};
% \end{tikzpicture}
\usepackage{tikz}

% Insertion d'images externes avec réglage de taille ou rotation.
% Avec pdfLaTeX, on utilise notamment des PDF, PNG ou JPEG.
%
% Exemple :
% \includegraphics[width=0.7\linewidth]{figures/grid.pdf}
%
% \linewidth désigne la largeur disponible à cet endroit.
\usepackage{graphicx}

% --------------------------------------------------
% Page layout and links
% --------------------------------------------------

% Ajoute « Annexe A », « Annexe B », etc. aux titres et au sommaire.
% Les sections d'annexe sont placées dans l'environnement appendices.
\usepackage[title,titletoc]{appendix}
\renewcommand{\appendixname}{Annexe}

% Définit les marges : ici 2,5 cm sur les quatre côtés.
\usepackage[margin=2.5cm]{geometry}

% Conserve des espacements verticaux naturels, même en recto verso.
% L'espace inutilisé reste en bas de page au lieu d'étirer le contenu.
\raggedbottom

% Numéros en pied de page, côté extérieur pour une reliure à gauche.
% LE : gauche des pages paires ; RO : droite des pages impaires.
\usepackage{fancyhdr}
\pagestyle{fancy}
\fancyhf{}
\fancyfoot[LE,RO]{\thepage}
\renewcommand{\headrulewidth}{0pt}
\renewcommand{\footrulewidth}{0pt}

% Même disposition pour les pages utilisant le style plain.
% Le style empty de la couverture et de la page blanche reste inchangé.
\fancypagestyle{plain}{%
    \fancyhf{}
    \fancyfoot[LE,RO]{\thepage}
    \renewcommand{\headrulewidth}{0pt}
    \renewcommand{\footrulewidth}{0pt}
}

% Gère les références stockées dans bibliography.bib.
% Les citations sont numérotées dans leur ordre d'apparition.
% Exemple dans le texte : \cite{stam1999stable}.
% Latexmk lance Biber pour construire la bibliographie.
\usepackage[backend=biber,style=numeric,sorting=none]{biblatex}
\addbibresource{bibliography.bib}

% Adapte les guillemets à la langue du document.
% Recommandé avec biblatex et babel.
\usepackage{csquotes}

% Rend cliquables les références et les entrées du sommaire.
% Fournit également les commandes \url et \href.
%
% hidelinks masque les couleurs et les cadres des liens,
% mais les liens restent cliquables.
% À charger généralement vers la fin du préambule.
%
% Exemple : \url{https://www.latex-project.org}
% Exemple : \href{https://www.latex-project.org}{LaTeX Project}
\usepackage[hidelinks]{hyperref}

% Met en gras italique les termes cliquables renvoyant à une explication.
% Premier argument : label cible ; second argument : texte affiché.
% Exemple : \conceptlink{annexes:forme_conservative}{forme conservative}
\newcommand{\conceptlink}[2]{%
    \hyperref[#1]{\textbf{\textit{#2}}}%
}
